\documentclass[a4paper]{article}

\usepackage{newtxtext}
\usepackage[utf8]{inputenc}
\usepackage[T1]{fontenc}
\usepackage{textcomp}
\usepackage{amsmath, amssymb}
\edef\restoreparindent{\parindent=\the\parindent\relax}
\DeclareSymbolFont{letters}{OML}{ztmcm}{m}{it}
\usepackage{parskip}
\restoreparindent
\newtheorem{problem}{Problem}
\newtheorem{example}{Example}
\newtheorem{lemma}{Lemma}
\newtheorem{theorem}{Theorem}
\newtheorem{problem}{Problem}
\newtheorem{example}{Example}
\newtheorem{definition}{Definition}
\newtheorem{lemma}{Lemma}
\newtheorem{theorem}{Theorem}

\begin{document}
    

Hola Pri,

Primero que nada te pido disculpas por no haberte respondido antes. Estaba de
viaje por trabajo, cuando volví tuve que acomodar muchas cosas, y quería tomarme
el tiempo de escribirte algo en detalle para corresponder a tu mensaje.

No sé si lo que tengo que decirte te va a alegrar o te va a dar ganas de tirarme
un zapatazo en la cabeza. Lo cierto es que desde mí subjetividad viví lo que
sucedió de una forma muy distinta a vos. Y si bien considero importante no negar
ni deslegitimar cómo fue tu experiencia personal, así como no dejar de hacerme
cargo de cómo contribuí a que fuera así, también me siento obligado a serte
honesto respecto a cómo yo lo viví.

Y lo cierto es que, para mí, más que como un «ghosteo» fue como un
distanciamiento que se dio más o menos espontáneamente, y que al menos yo
sospeché que iba de ambos lados. Me fui de viaje a Corrientes y en ese viaje no
hablamos casi nada (o nada, no recuerdo bien); a mi regreso ninguno activó, y
justo el día que volví o el siguiente nos cruzamos en ciudad universitaria y,
aunque nos saludamos con la mejor onda, tampoco surgió nada después. Yo en lo
personal sentí que fue algo que se fue dando naturalmente; más que como «le dejé
de hablar» lo viví como «nos dejamos de hablar». El por qué de ese
distanciamiento, que al menos desde mi experiencia se sintió como algo mutuo, no
lo sé; sentí que fue algo que simplemente a veces pasa: empezar a salir con
alguien y, con el paso de las semanas, sentir que se hablan menos y menos hasta
que «queda ahí».

Ahora, quiero aclarar una cosa: en ningún momento sentí nada negativo respecto a
vos. Es más, hasta el día de hoy te recuerdo como una persona cálida, linda,
divertida, y \textit{muy} teatrera en el mejor sentido del término. Si te
hubiera cruzado en la calle, sin imaginarme que vos habías vivido lo que pasó
como un ghosteo, te hubiera saludado con todo el cariño del mundo. Creo que un
poco por eso también te seguí en Instagram y te likeé muchas historias, y ahora
me siento un bobo porque yo lo hacía pensando que estaba todo bien, que
simplemente las cosas se habían apagado pero sin malos sentimientos.

No quiero que esta carta se interprete como un «te estás comiendo un viaje, lo
que decís no sucedió». Si vos lo sentiste como un ghosteo yo tengo que
replantearme qué pude haber hecho distinto para que no lo sintieras así, y en
qué me confundí cuando lo interpreté como algo diferente. Te quiero pedir perdón
porque no quise que te sintieras mal y de haber interpretado correctamente la
situación habría actuado de otro modo. 

De mi parte está todo bien, te tengo aprecio y te considero una persona hermosa.
Si de tu parte no está todo bien, lo respeto y lo entiendo. Sólo espero que
entiendas que nada de esto fue un acto deliberado para hacerte sentir mal sino
simplemente dos personas que interpretaron de manera muy diferente lo que estaba
sucediendo.

---Santi




















































\end{document}
